\Def Системой непересекающихся множеств назовём структуру данных, которая в любой момент поддерживает некоторый набор множеств $\mathcal{A} = \{A_1, A_2, \dots A_n\}: i \neq j \Rightarrow A_i \cap A_j = \varnothing$, а также позволяет выполнять три операции:
\begin{itemize}
    \item $\texttt{MAKE\_SET(}x\texttt{)}$ -- создает множество с единственным элементом $x$, который не принадлежит ни одному множеству.
    \item $\texttt{SAME(}x, y\texttt{)}$ -- возвращает \texttt{true}, если $\exists i: x, y \in A_i$ (то есть два элемента лежат в одном и том же множестве).
    \item $\texttt{UNITE(}x, y\texttt{)}$ -- выполняет объединение двух множеств, в которых лежат $x$ и $y$.
\end{itemize}
Элементы множества $A_i$ иногда назвают \textit{представителями множества}.

\Exe Предложите алгоритм, который выделяет компоненты связности в графе, с помощью СНМ.

\section{Реализация с помощью леса непересекающихся множеств}

Далее рассмотрим одну из реализаций СНМ -- \textit{лес непересекающихся множеств}.

Каждое множество будем представлять в виде дерева, которое состоит из элементов этого множества.

\begin{figure}[H]
    \centering
    \begin{tikzpicture}[
        node distance=1cm,
        state/.style={circle,draw,fill=white!20,minimum size=0.5cm}
    ]

    \node[state] (c) {$c$};
    \node[state] (h) [below left=of c] {$h$};
    \node[state] (e) [below right=of c] {$e$};
    \node[state] (b) [below=of h] {$b$};

    \draw[->,thick] (c) edge[loop above] (c);
    \draw[->,thick] (h) -- (c);
    \draw[->,thick] (e) -- (c);
    \draw[->,thick] (b) -- (h);


    \node[state] (f) [right=4cm of c] {$f$};
    \node[state] (d) [below=of f] {$d$};
    \node[state] (g) [below=of d] {$g$};

    \draw[->,thick] (f) edge[loop above] (f);
    \draw[->,thick] (d) -- (f);
    \draw[->,thick] (g) -- (d);

    \end{tikzpicture}
    \caption{Пример СНМ для $\mathcal{A} = \{ \{c, h, b, e\}, \{f, d, g\}\}$}
\end{figure}

Тогда в такой парадигме интерфейс реализуется следующим образом:
\begin{itemize}
    \item $\texttt{MAKE\_SET(}x\texttt{)}$ -- создание узла.
    \item $\texttt{SAME(}x, y\texttt{)}$ -- для $x$ и $y$ находим корни деревьев, в которых они лежат, $r_x$ и $r_y$ при помощи поднятия вверх. Затем сравниваем эти корни, если $x$ и $y$ лежат в одном множестве, то $r_x = r_y$, иначе $r_x \neq r_y$.
    \item $\texttt{UNITE(}x, y\texttt{)}$ -- также как в \texttt{SAME} находим корни $r_x$ и $r_y$, а затем подвешиваем $r_y$ в качестве ребенка к $r_x$.
\end{itemize}

Итак, исследуем время работы нашей структуры данных. Очевидно, \texttt{MAKE\_SET} работает за $\O(1)$.
Но вот с \texttt{UNIT} и \texttt{SAME} уже не всё так просто, потому что время их работы зависит напрямую от высот деревьев, в которых лежат $x$ и $y$.
И не сложно заметить, если взять два бамбука высоты $h$ и вызвать \texttt{SAME} от самых глубоких вершин, то время работы будет $\Theta(h)$.

На самом деле можно ускорить работу этих операций вплоть до <<почти константы>>. Для этого надо использовать две эвристики: \textit{ранговую} и \textit{сжатие путей}.
